% 2026-09-23 — order-1296 complex-reflection no-go.
\subsection*{The extended qutrit \(1296\)-group is not Shephard--Todd \(G_{26}\)}

The order alone suggests a tempting comparison with the three-dimensional
Hessian complex-reflection group \(G_{26}\), also of order \(1296\).  The
comparison is now ruled out by an exact group invariant.

The retained-phase extended qutrit Clifford group certified above is
\[
G_{\rm ext}\cong H_{27}:GL(2,3),
\qquad |G_{\rm ext}|=1296,
\]
with
\[
\boxed{Z(G_{\rm ext})=1.}
\]
Its determinant-minus-one coset inverts the normal Pauli phase subgroup
\(C_3\), so that subgroup is no longer central in the full extension.

By contrast, the standard Shephard--Todd classification gives
\[
\boxed{
G_{26}\cong C_2\times\bigl(3^{1+2}\!\cdot SL(2,3)\bigr),
}
\]
again of order \(1296\).  In particular \(G_{26}\) contains an explicit
central involution. Therefore
\[
\boxed{
G_{\rm ext}\not\cong G_{26}.
}
\]
The natural quotient pictures also separate: the qutrit group's normal
phase-\(C_3\) quotient is \(AGL(2,3)\) of order \(432\), whereas the scalar
collineation quotient of \(G_{26}\) is the classical Hessian group of order
\(216\).

\paragraph{Boundary.}
This no-go concerns the full \(1296\)-element anti-linear extension only.
It does not rule out the separately known Hessian/\(G_{25}\) relationship of
the determinant-one \(648\)-element core.  Equal orders are not being used
as identifications in either direction.
